% !TEX root = ../../../tree_notes.tex
% =====================================================================
%  12.3  The Centroid: the Tree's Centre of Mass   sec:c09-centroid
%  Exports: def:c09-centroid  thm:c09-centexist  thm:c09-centfind
%           lst:c09-cent  fig:c09-centcut  ex:c09-cent-run
%           prob:c09-cent  lst:c09-2079-soln   (practice: CF 342E, preview)
%  This section ENDS at "find one centroid + why it is safe to recurse on".
%  The decomposition itself (recursing) is Chapter 24's (chref{c18}) and
%  is previewed in prose only -- no listing, no template.
%  Hand-check bonus: the running tree T has TWO adjacent centroids {1,3};
%  the section uses that live, including the <= vs < boundary trap.
% =====================================================================
\section{The Centroid --- Cut at the Balance}\label{sec:c09-centroid}

\begin{previouslybox}
\begin{enumerate}[nosep, leftmargin=1.6em]
  \item \coderef{c09-dfs3} produced $\size(u)$ for every $u$ under a chosen root. Once
        $u$ is \emph{deleted}, which of its neighbours' sides hold $\size(c)$ nodes and
        which hold $n-\size(u)$? (children: their $\size$; the parent side: the complement)
  \item A path of $4$ nodes: cutting at either middle node leaves pieces of size $1{+}2$;
        is the largest piece $2 \le n/2$? (yes --- both middles qualify; hold that thought)
  \item What does a divide-and-conquer algorithm need from its split to guarantee
        $\bigO{\log n}$ levels? (each part at most a fixed fraction of the whole)
  \item Could a tree have \emph{no} node whose removal caps every piece at $n/2$? What
        would a proof of the opposite look like? (walk toward the heavy piece; sizes
        shrink --- this section's \thmref{c09-centexist})
\end{enumerate}
\end{previouslybox}

\begin{topicintro}
The diameter measured the tree's \emph{extremes}; the centroid is about its
\emph{middle} --- but a middle defined by counts, not distances. A node $c$ is a
\term{centroid} if deleting it leaves no connected piece larger than $n/2$
(\defref{c09-centroid}): the point where the tree's mass balances. Two facts make it the
most load-bearing definition in contest trees. First, \emph{existence}: every tree has
one, found by walking toward whatever side is heavy (\thmref{c09-centexist}) --- and at
most two, which then sit adjacent. Second, \emph{utility}: a centroid is exactly the
split a divide-and-conquer needs; recursion that removes a centroid halves the problem
at every level, which is why \thmref{c09-centexist}'s proof doubles as the blueprint for
centroid \emph{decomposition} --- a technique this chapter hands to Chapter~24
(\chref{c18}) after one preview paragraph.

Everything computational here is $\size[\cdot]$ re-read from a different angle: the
pieces created by deleting $u$ are its child subtrees and one \emph{upward} remainder
$n-\size(u)$ --- forgetting that remainder is this section's signature bug (Trap~1). We
close on CSES~2079 (\probref{c09-cent}), whose whole input is ``print a centroid,'' and
the running tree $T$ supplies a gift: its centroid is \emph{not unique} --- $\{1,3\}$
is an adjacent pair, each with a piece of \emph{exactly} $n/2$ --- so the boundary
between ``$\le n/2$'' and ``$< n/2$'' is not pedantry, it is the difference between
finding and missing the answer.
\textbf{Prerequisites.} \secref{c09-measures} ($\size$ and the parent-guard DFS; the
\coderef{c09-iter} discipline); nothing else --- no distances anywhere in this section.

\textbf{After this section you can}
\begin{itemize}[nosep]
  \item list the pieces created by deleting a node and test centroid-ness in $\bigO{\deg}$; \hfill\lv{3}
  \item prove existence (walk-to-heavy) and the 1-or-2-adjacent rule; \hfill\lv{4}
  \item find a centroid in one DFS via \code{maxpart}, and state why recursion on
        centroid pieces halves the depth; \hfill\lv{3}--\lv{4}
\end{itemize}
\end{topicintro}

\subsection{Intuition}

Hold the tree by one node and shake it: everything that stays in your hand is the node;
everything that falls is a \emph{piece}. The centroid is the grip where no falling strand
is heavier than half the tree. That image fixes the bookkeeping: with a root chosen
arbitrarily, the pieces falling from $u$ are precisely the child-subtrees of $u$ (sizes
already computed) plus one strand going \emph{up} through the parent, containing
$u$'s ancestors and everything not in $u$'s subtree --- $n - \size(u)$ nodes. Three
quantities to combine, all from the same $\size$ table.

Existence is a descent, not a search: if the grip at $u$ fails, some piece $P$ has more
than $n/2$ nodes. Grip any node inside $P$ instead --- specifically $u$'s neighbour on
the path into $P$ --- and you have moved \emph{into} the majority, so the piece you
came from is now minority ($< n/2$) and the only possible failure is deeper inside $P$.
Trees are finite, the majority region strictly shrinks, the walk stops: at rest every
piece is a minority, which is exactly the definition. Two grips can both rest only when
the tree divides $n/2 {+} n/2$ across the edge between them --- the pair case, with $T$
as its exhibition model.

Why should a problem about \emph{counting} care about \emph{halving}? Because half is
the divide in divide-and-conquer: solve anything by ``answer for paths through the
grip-node, then recurse on the pieces,'' and \thmref{c09-centexist} is the guarantee
that the recursion is $\bigO{\log n}$ deep instead of $n$ --- the difference between a
$\bigO{n \log n}$ algorithm and an $\bigO{n^2}$ one on a chain.

\subsection{Definition and the Two Structure Theorems}

\begin{definition}{Centroid}{c09-centroid}
A node $c$ of a tree $T$ on $n$ nodes is a \term{centroid} if every connected component
of $T - c$ has at most $\lfloor n/2 \rfloor$ nodes. Equivalently, rooting $T$ anywhere
and writing, for each node $u$,
\[
  \mathit{maxpart}(u) \;=\; \max\!\Big(\, \underbrace{n - \size(u)}_{\text{upward piece}},\;\;
        \max_{c \in \children(u)} \size(c) \,\Big),
\]
with $\mathit{maxpart}(u) := 0$ when $u$ is the whole tree ($n=1$), $c$ is a centroid
iff $\mathit{maxpart}(c) \le \lfloor n/2 \rfloor$.
\end{definition}

\noindent\textbf{In plain words.} Delete $c$, count the shards, the biggest must be at
most half. The formula is the shard list of a grip at $u$ read off the $\size$ table:
each child brings its subtree, the parent brings \emph{the rest} --- and $\sum$ of all
pieces is $n-1$, so at most one piece can exceed $n/2$ (two would need $> n$ fingers).
That single-majority fact powers every proof below.

\begin{theorem}{Existence, and at most two (adjacent)}{c09-centexist}
Every tree has a centroid. The set of centroids has size one or two; when two, they are
adjacent, $n$ is even, and deleting either leaves a piece of exactly $n/2$ nodes.
\end{theorem}
\begin{proofbox}
\emph{Existence (walk-to-heavy).} Start at any node. If some piece at $u$ holds $> n/2$
nodes, cross the edge into that piece (its other end is a neighbour $v$, and at most one
such neighbour exists by the majority count above). Moving from $u$ to $v$: the piece
that was ``the side of $v$'' at $u$ becomes the upward piece at $v$ with the \emph{same}
size, while every other piece at $v$ sits strictly inside it. The majority-piece size
therefore strictly decreases at every step (from $> n/2$ to its sub-piece), so by
well-foundedness of the integers the walk halts --- halting means no piece exceeds
$n/2$, i.e.\ a centroid.

\emph{At most two, adjacent.} Suppose centroids $u \ne v$ and let $A$ = the component of
$T-u$ containing $v$, $B$ = the component of $T-v$ containing $u$. Both paths' interiors
lie in $A \cap B$, and $A \cup B = V$, so $n = |A| + |B| - |A \cap B|$. Centroid-ness
gives $|A| \le n/2$ and $|B| \le n/2$, hence $|A \cap B| \le |A| + |B| - n \le 0$: the
intersection is empty, which happens exactly when $u,v$ are \emph{adjacent} (no interior
nodes; then $A, B$ partition $V$ and $|A| = |B| = n/2$, forcing even $n$). Three or more
centroids would make some pair non-adjacent, contradicting the adjacency conclusion.
\end{proofbox}

\begin{theorem}{One DFS finds a centroid}{c09-centfind}
Run any $\size$ pass (\coderef{c09-dfs3} or \coderef{c09-iter}); then $u$ is a centroid
iff $\mathit{maxpart}(u) \le \lfloor n/2\rfloor$ with the $\bigO{\deg}$ formula of
\defref{c09-centroid}, computable inside the same DFS --- total $\Theta(n)$.
Equivalently, start at the root and repeatedly descend into the (necessarily unique)
child with $\size > n/2$; the terminal node is a centroid.
\end{theorem}
\begin{proofbox}
The iff is \defref{c09-centroid} restated; the pieces at $u$ are the child subtrees
($\size$ known) and the upward complement $n - \size(u)$ ($\size$ known), so one max
over $\deg(u)$ values decides $u$, and one scan decides all. For the descent form: at
the root, crossing into child $v$ with $\size(v) > n/2$ is exactly a walk-to-heavy step:
the upward piece at $v$ is $n - \size(v) < n/2$, and every child piece of $v$ sits
inside $v$'s subtree; so the descent \emph{is}
the walk of \thmref{c09-centexist}'s first proof, restricted to the rooted orientation
(the heavy piece, if not a child, is the upward one, which the descent never creates:
we only enter subtrees of size $> n/2$, whose complements are automatically $\le n/2$).
Terminating at a node with no heavy child then means all pieces $\le n/2$.
\end{proofbox}

\begin{lstlisting}[caption={maxpart in one DFS: the upward piece n - sz[u] is the line people forget}, label={lst:c09-cent}]
// sz[] from dfs3; then a second pass (or same pass, after child returns):
vector<int> szc, maxpart;
void dfs_cent(int u, int p) {
    szc[u] = 1;  maxpart[u] = 0;
    for (int v : adj[u]) if (v != p) {
        dfs_cent(v, u);
        szc[u]    += szc[v];
        maxpart[u] = max(maxpart[u], szc[v]);     // child pieces
    }
    maxpart[u] = max(maxpart[u], n - szc[u]);     // THE UPWARD PIECE  <-- Trap 1
}
// centroid test: 2 * maxpart[u] <= n   (integer-safe form of maxpart <= n/2)
// call dfs_cent(r, -1) for ANY r; maxpart is a property of u, not of the root.
\end{lstlisting}

\begin{insightbox}[Centroid vs centre: same word, different mathematics]
The \term{centre} (\secref{c09-radius}) minimises \emph{distance to the farthest node};
the \term{centroid} minimises \emph{the largest remaining piece}. They disagree constantly:
take a path $1-2-3-4-5$ and hang $20$ leaves off node~$2$ ($n = 26$). The centroid is $2$:
deleting it leaves $20$ singleton leaf-pieces and the path shards $\{1\}$, $\{3,4,5\}$
--- biggest piece $3 \le 13$ --- while deleting any path node $3$ or $4$ strands the
$22$-node blob containing the star ($> 13$). The centre is $3$: its eccentricity is
$2$ (the leaves are $2$ steps away, the path ends $2$ too), whereas node~$2$ sits $3$
from node~$5$. Coincidence looks universal only when
examples are small: $T$'s centre pair $\{1,3\}$ equals its centroid pair --- prove it is
luck by finding a tree where they differ, e.g.\ the broom.
\end{insightbox}

\subsection{Why Half Is the Magic Number (decomposition preview --- Chapter 24's)}

Suppose you must process ``all paths of $T$ satisfying property $P$.'' Grip a centroid
$c$: every path either passes through $c$ ( enumerable with $c$'s piece-structures in
linear-ish time) or lives entirely inside one piece --- and every piece has
$\le n/2$ nodes (\thmref{c09-centexist}), so recursing \emph{on pieces} halves the input
at every level: depth $\le \log_2 n$, total work $\bigO{n \log n}$ for any per-level
linear schedule. That sentence is the entire \term{centroid decomposition} in embryo ---
the data structure built on it, the recursion bookkeeping, and the path-counting
applications belong to Chapter~24 (\chref{c18}), and this chapter deliberately stops
here: \textbf{one centroid, found in }$\Theta(n)$\textbf{, is all you need to know to
finish this book's early chapters}. (The only preview exercise we assign, CF~342E in
the practice set, is explicitly gated on that chapter.)

\subsection{Worked Example}

\begin{example}{Centroids of $T$: the two-adjacent case on the board}{c09-cent-run}
$n = 10$, $\lfloor n/2 \rfloor = 5$, $\size[\cdot]$ from \figref{c09-triples}. Read
$\mathit{maxpart}$ off the definition at each candidate (root $1$):
\[
\begin{array}{c|l|c|c}
u & \text{pieces at } u & \mathit{maxpart} & \text{centroid?} \\
\hline
1 & \size(2){=}4,\; \size(3){=}5,\; \text{up }{=}0 & 5 & \le 5\ \checkmark \\
3 & \size(6){=}1,\; \size(7){=}3,\; \text{up }{=}10{-}5{=}5 & 5 & \le 5\ \checkmark \\
2 & \size(4){=}1,\; \size(5){=}2,\; \text{up }{=}10{-}4{=}6 & 6 & \times \\
7 & \size(9){=}2,\; \text{up }{=}10{-}3{=}7 & 7 & \times \\
9 & \size(10){=}1,\; \text{up }{=}10{-}2{=}8 & 8 & \times
\end{array}
\]
Centroids $=\{1,3\}$ --- adjacent, $n$ even, and indeed deleting $1$ leaves the piece
$\{3,6,7,9,10\}$ of exactly $5 = n/2$, deleting $3$ leaves $\{1,2,4,5,8\}$ of exactly
$5$: the pair case of \thmref{c09-centexist} with every clause exercised. A solver
testing ``$<$ instead of $\le$'' finds \emph{no} centroid on $T$; a solver using the
\emph{strict majority} descent stops at $1$ or $3$ and must still accept. \end{example}

\begin{figure}[H]
\centering
\begin{tikzpicture}[diagram, scale=0.82, transform shape,
    tnode/.style={circle, draw=primaryblue, line width=0.9pt, fill=primaryblue!8,
                  minimum size=6.2mm, font=\small\bfseries, text=primaryblue!60!black},
    cen/.style={tnode, draw=softpurple, fill=softpurple!22, line width=1.2pt,
                text=softpurple!60!black}]
  \node[cen] (n1) at (0,0) {$1$};
  \node[tnode] (n2) at (-3,-1.6) {$2$};
  \node[cen] (n3) at ( 3,-1.6) {$3$};
  \node[tnode] (n4) at (-5.2,-3.2) {$4$};
  \node[tnode] (n5) at (-1.6,-3.2) {$5$};
  \node[tnode] (n6) at ( 1.6,-3.2) {$6$};
  \node[tnode] (n7) at ( 5.2,-3.2) {$7$};
  \node[tnode] (n8) at (-1.6,-4.8) {$8$};
  \node[tnode] (n9) at ( 5.2,-4.8) {$9$};
  \node[tnode] (n10) at ( 5.2,-6.4) {$10$};
  \draw[tedge] (n1)--(n2) (n1)--(n3) (n2)--(n4) (n2)--(n5)
               (n3)--(n6) (n3)--(n7) (n5)--(n8) (n7)--(n9) (n9)--(n10);
  \draw[guide] (n1)--(n3);
  \node[annot, font=\scriptsize, text=softpurple!70!black, anchor=south east]
    at (n1.north west) {cut here: pieces $4 \mid 5$};
  \node[annot, font=\scriptsize, text=softpurple!70!black, anchor=south west]
    at (n3.north east) {or here: $5 \mid 3 \mid 1$};
  \node[annot, font=\scriptsize, anchor=north, text=darkgray]
    at (1.5,-2.15) {the centroid edge: the $5{+}5$ split};
\end{tikzpicture}
\caption*{Both centroids of $T$ in purple; the dashed guide marks the edge whose removal
realises the $n/2 + n/2$ split that \thmref{c09-centexist} says must exist whenever the
centroid is not unique. Cutting at $1$: pieces of $4$ and $5$ nodes; cutting at $3$:
pieces of $5$, $3$, $1$. Neither exceeds $\lfloor 10/2 \rfloor = 5$.}
\label{fig:c09-centcut}
\end{figure}

\begin{tracebox}[Walk-to-heavy from node 8 (the descent proof, live)]
at 8: pieces \{5-side: 9 nodes\} -- 9 > 5, cross to 5.\\
at 5: pieces \{up-side: 8\}, \{8: 1\} -- 8 > 5, cross to 2.\\
at 2: pieces \{up: 6\}, \{4: 1\}, \{5-side: 2\} -- 6 > 5, cross to 1.\\
at 1: pieces \{4-side: 4\}, \{3-side: 5\} -- none > 5. STOP: centroid 1.
Sizes followed: 9, 8, 6, then rest --- strictly decreasing majorities, exactly the
proof's well-founded measure. Crossing one more step to 3 would show its piece of 5 is
merely \emph{equal} to $n/2$: legal, hence the second centroid. \checkmark
\end{tracebox}

\subsection{Variations and Twists}

\begin{itemize}[nosep]
  \item \textbf{The descent as code:} replace \coderef{c09-cent}'s scan by ``while a
        child has $\size > n/2$, descend into it'' --- $\bigO{n}$ after $\size$, with
        $\bigO{h}$ no extra pass; the two centroids reveal themselves as the pair across
        a $\size = n/2$ edge. \hfill\lv{3}
  \item \textbf{Node-weighted centroid:} give nodes weights, cap every piece at $W/2$
        ($W = \sum$ weights); walk-to-heavy is word-for-word unchanged (the majority
        count is now by weight). The trie-with-counts variants of Chapter~11 use
        exactly this. \hfill\lv{4}
  \item \textbf{Edge variant:} ``remove an \emph{edge} to balance'' uses the same
        $\size$ table ($n - 2\size(v)$ drives the imbalance) but \emph{no} existence
        theorem --- a balanced edge need not exist; don't import centroid guarantees. \hfill\lv{4}
  \item \textbf{Centroid of a binary tree is not the ``middle'':} on the broom of the
        insight box the centroid sits at the star's hub; if your mental image says
        ``geometric middle,'' you are thinking of the centre --- next section. \hfill\lv{3}
\end{itemize}

\subsection{Common Pitfalls}

\begin{warningbox}[Trap 1 -- forgetting the upward piece]
$\mathit{maxpart}(u) = \max$ over \emph{children} only: on $T$ rooted at~$1$, node~$3$
looks like it has max piece $3$ (child $7$), and node~$2$ looks fine at~$2$ --- while
the real values are $5$ and $6$. The parent side is a piece the moment you delete $u$:
$T - u$ does not know where the root was. Symptom: your code returns a leaf-adjacent
node on chains (where $\size$ is tiny down-root but $n - \size$ towers behind). The
line marked \emph{Trap~1} in \coderef{c09-cent} is the whole fix.
\end{warningbox}

\begin{warningbox}[Trap 2 -- strict $<$ where the definition says $\le$]
``Every piece has size $\le n/2$'' includes equality --- and equality is \emph{exactly}
where the two-centroid case lives ($T$'s nodes~$1,3$, both with a piece of $5$). Write
\code{2*maxpart[u] <= n} (integer arithmetic: no float, no rounding debate) and
\texttt{<} becomes unmissable. A ``$<$'' solver prints nothing on $T$ and crashes on
the judge's ``print any centroid'' the way silence always does.
\end{warningbox}

\begin{warningbox}[Trap 3 -- claiming uniqueness, or claiming it never happens]
Both overreaches break proofs of \emph{later} steps: algorithms that ``the'' centroid
refers to must handle the pair (pick either --- deleting either bounds pieces by $n/2$,
which is all any recursion needs), and the adjacency clause is what bounds the \emph{second}
piece-structure when both are removed. Remember \thmref{c09-centexist} as: ``one or two,
adjacent, split $n/2{+}n/2$ iff two.'' \hfill($\le n/2$ boundary cases patrol the WA
clinic at the chapter end.)
\end{warningbox}

\subsection{Problem Lab}

The anchor is CSES~2079, \probref{c09-cent}: read a tree, print one centroid --- the
definition with zero garnish. The garnish is entirely in what you must not forget:
the upward piece (Trap~1), the boundary (Trap~2), and the stack discipline for
$n = 2\cdot 10^{5}$ chains (\coderef{c09-iter}, already sworn in at 12.1).

% ---------------------------------------------------------------------------
\solvedproblem{c09-cent}{Finding a Centroid}{CSES 2079}
  {https://cses.fi/problemset/task/2079}{CSES · L3}
  {subtree sizes, maxpart with upward piece, iterative DFS discipline}

\begin{problemstatement}[Finding a Centroid (statement and sample verified against the live problem)]
Given a tree of $n$ nodes ($1 \le n \le 2\cdot 10^5$), print \emph{any} centroid: a node
whose removal leaves every component with at most $n/2$ nodes (\defref{c09-centroid}).
\textbf{Sample}: $n = 5$, edges $1\!-\!2$, $2\!-\!3$, $3\!-\!4$, $3\!-\!5$
$\Rightarrow$ output \texttt{3} (removing~$3$: pieces $\{1,2\}$, $\{4\}$, $\{5\}$,
sizes $2,1,1 \le \lfloor 5/2 \rfloor$; the centroid here is unique).
\end{problemstatement}

\begin{attemptgate}[Attempt it yourself first -- \lv{3}, give yourself 15 minutes]
One DFS already handed you $\size[\cdot]$; the rest is reading the shard list correctly.
\peek{After deleting a non-root node $u$, how many pieces exist and what are their sizes
in terms of $\size$?}
\nudge{One per child ($\size(c)$) plus one upward piece of $n - \size(u)$. Take the max
over all --- that's $\mathit{maxpart}$ of \defref{c09-centroid}.}
\fullhint{Root at~1, fill $\size$ with \coderef{c09-dfs3} (or \coderef{c09-iter}:
$n=2\cdot10^5$ and CSES happily serves chains), compute $\mathit{maxpart}$ per
\coderef{c09-cent}, print any $u$ with \code{2*maxpart[u] <= n}. The sample: $\size$ at
$3$ is $3$, so up-piece $=2$; child pieces $1,1$; $\mathit{maxpart}(3)=2$,
$2\cdot 2 = 4 \le 5$ ✓; check $2$ fails ($\mathit{maxpart}=3$, the $3$-node side below).}
\end{attemptgate}

\solstep{Understanding the Problem}
Output is one node label; \thmref{c09-centexist} guarantees at least one exists, and
accepting \emph{any} means the pair case ($\{1,3\}$ on $T$) can print either. No
distance is computed anywhere --- if your plan mentions ``farthest,'' you have switched
to \secref{c09-radius}'s problem. The real engineering: $n \le 2\cdot 10^5$ with chains
legal, so budget the stack (\coderef{c09-iter}) before the first submit.

\solstep{Key Observations}
\begin{observation}
The pieces of $T-u$ under a fixed root are exactly: one per child ($\size(c)$) and the
upward remainder ($n - \size(u)$) --- so centroid-ness is a $\bigO{\deg}$ readoff of the
$\size$ table (\defref{c09-centroid}), $\Theta(n)$ for all nodes (\thmref{c09-centfind}).
\end{observation}
\begin{observation}
Root-choice independence: $\mathit{maxpart}(u)$ enumerates the same \emph{physical}
pieces whichever root was used --- the root is bookkeeping, not part of the question.
\end{observation}
\begin{observation}
$2 \cdot \mathit{maxpart} \le n$ keeps the boundary in integers; equality is legal and
marks the two-centroid edge when it holds for two nodes.
\end{observation}

\solstep{Think It Through With Me}
\begin{thinkbox}
\deadend{Try every node, deleting and re-counting components by BFS per candidate.}
{$\bigO{n^2}$: $2\cdot 10^5$ nodes $\Rightarrow \sim 4\cdot 10^{10}$ --- TLE; it also
re-derives $\size$ $n$ times instead of once.}
\deadend{Binary-search the ``balance point'' along the diameter path.}
{Centroids need not lie on the diameter you found in 12.2 (check: the broom's centroid
is the star hub, distance-geometry plays no role), and monotonicity along the path is
not guaranteed for off-path mass --- you would be searching the wrong line.}
\aha{One $\size$ pass + the upward-piece formula: $\mathit{maxpart}$ in the same DFS
(\coderef{c09-cent}), first node with $2\mathit{maxpart} \le n$ wins --- $\Theta(n)$,
and \thmref{c09-centexist} says the winner exists so the loop never prints nothing.}
\end{thinkbox}

\solstep{Solution}
Iterative two-phase DFS (\coderef{c09-iter}'s machinery with a centroid payload):
preorder plants parent/depth, reversed replay fills $\size$ and, per node,
$\mathit{maxpart} = \max(\text{child pieces}, n - \size)$; scan and print. Correctness:
\thmref{c09-centfind} (iff), \thmref{c09-centexist} (non-empty answer set).

\solstep{Proof of Correctness}
\begin{proofbox}[The printed node is a centroid]
For the printed $u$ the code evaluated $\mathit{maxpart}(u)$ as the max of $\size(c)$
over $u$'s children and $n-\size(u)$. These are exactly the sizes of the connected
components of $T - u$: each child subtree survives as a component of its recorded size,
and the remaining nodes (ancestors and all other branches) form one connected piece of
$n - \size(u)$ (connectedness: paths avoiding $u$ route through the parent side).
Acceptance of $u$ requires $2\,\mathit{maxpart}(u) \le n$, i.e.\ every component
$\le \lfloor n/2\rfloor$, which is \defref{c09-centroid}'s condition verbatim. Such a
$u$ exists by \thmref{c09-centexist}, so the scan's ``found'' branch always fires; the
$\size$ values themselves are correct by \thmref{c09-onepass}'s induction (the reversed
replay is a legal fold order, \secref{c09-measures}).
\end{proofbox}

\solstep{Implementation}
\begin{lstlisting}[caption={CSES 2079 -- iterative size pass, then the maxpart scan}, label={lst:c09-2079-soln}]
#include <bits/stdc++.h>
using namespace std;
int main() {
    int n;  scanf("%d", &n);
    vector<vector<int>> adj(n);
    for (int i = 0; i < n - 1; i++) {
        int a, b;  scanf("%d %d", &a, &b);  --a; --b;
        adj[a].push_back(b);  adj[b].push_back(a);
    }
    if (n == 1) { printf("1\n"); return 0; }      // the trivial centroid
    vector<int> par(n, -1), sz(n), order;
    order.reserve(n);
    vector<int> st{0};
    while (!st.empty()) {                          // preorder (cf. lst:c09-iter)
        int u = st.back();  st.pop_back();
        order.push_back(u);
        for (int v : adj[u]) if (v != par[u]) { par[v] = u; st.push_back(v); }
    }
    int best = -1;
    for (int i = (int)order.size() - 1; i >= 0; --i) {   // children first
        int u = order[i];
        sz[u] = 1;  int mx = 0;
        for (int v : adj[u]) if (v != par[u]) {
            sz[u] += sz[v];
            mx = max(mx, sz[v]);                   // child pieces
        }
        mx = max(mx, n - sz[u]);                   // the upward piece -- Trap 1
        if (2 * mx <= n) best = u;                 // any centroid; keeps last found
    }
    printf("%d\n", best + 1);
}
\end{lstlisting}

\solstep{Trace}
\begin{tracebox}[Official sample: 5 nodes, edges 1-2, 2-3, 3-4, 3-5]
Root 1 (label 0-based node 0). preorder e.g. 0,1,2,4,3; replay backwards:
node 4 (label 5): sz=1, mx = max(up 5-1=4) = 4 -> 8<=5? no.
node 3 (label 4): sz=1, mx=4 -> no.
node 2 (label 3): children 4,5: sz=3, mx = max(1,1, up 5-3=2) = 2 -> 4<=5 YES, best=2.
node 1 (label 2): child 3: sz=4, mx = max(3, up 1) = 3 -> 6<=5 no.
node 0 (label 1): child 2: sz=5, mx = max(4, up 0) = 4 -> no.
print best+1 = \textbf{3} \checkmark\ matches the official output (unique centroid;
compare \exref{c09-cent-run} where the scan would find two).
\end{tracebox}

\solstep{Complexity}
\begin{complexitybox}
Time: one preorder pass + one replay + per-node $\bigO{\deg}$ scans $= \Theta(n)$
(\thmref{c09-centfind}). Memory: $\Theta(n)$ arrays, $\bigO{1}$ call stack by design
(explicit stack --- the chain case that kills recursive submissions on CSES is a legal
input here). Verdict: \ac. Note what \emph{wasn't} needed: distances, diameter, any
tie-breaking wisdom about which of two centroids to print.
\end{complexitybox}

\solstep{Pitfalls and Variants}
\begin{warningbox}
Live-fire recap for the traps: forgetting \code{n - sz[u]} makes the chain's
$\mathit{maxpart}$ collapse to child-side values and the program prints a middle-ish
\emph{child} rather than failing loudly; \code{2*mx < n} prints nothing on even-split
cases like $T$. Variants: print \emph{all} centroids (collect the found set --- its size
is $1$ or $2$, \thmref{c09-centexist}, a good assert); node-weighted version
(\emph{Variations}); and CSES~1130-style guarantees about the pieces, which the mixed
set revisits.
\end{warningbox}

\subsection{Practice Problems}

\begin{stuckbox}
\begin{itemize}[nosep, leftmargin=1.4em]
  \item Does your $\mathit{maxpart}$ include $n - \size(u)$? Test on node~$2$ of $T$
        (should be $6$, not $2$).
  \item Is your acceptance test $\le$ in integer form \code{2*mx <= n}?
  \item Does your DFS survive a $2\cdot 10^{5}$ chain without recursion?
\end{itemize}
\end{stuckbox}

\begin{practiceset}
  \practice{Xenia and Tree}{Codeforces 342E}
    {https://codeforces.com/problemset/problem/342/E}{CF 2400 · L6}
    {Paint nodes red one at a time; each query asks the minimum distance from a given
    node to any red node. \textbf{Gate warning:} the clean solution is centroid
    decomposition --- Chapter~24 (\chref{c18}) --- and you should \emph{not} grind this
    before that chapter; it is listed as the section's promise-of-a-problem, the thing
    ``find one centroid'' becomes once recursion on the pieces is automated. When you
    do return: store, at every centroid ancestor of a node, the distance to the nearest
    red in its piece; the centroid bound of \thmref{c09-centexist} is what makes the
    ancestor chain $\bigO{\log n}$ long. A pre-decomposition taste: answer all queries
    \emph{offline} by multi-source BFS per batch of $\bigO{\sqrt{n}}$ repaints ---
    sqrt-decomposition is the honest 12.3-level stretch.}
\end{practiceset}

\begin{recallbox}[10-second refresher -- 12.3 The Centroid]
$c$ is a centroid if every piece of $T-c$ has $\le \lfloor n/2\rfloor$ nodes
(\defref{c09-centroid}): pieces $=$ child subtrees $+$ the upward $n-\size(u)$. Every
tree has one (walk-to-heavy, \thmref{c09-centexist}); at most two, adjacent, $n/2{+}n/2$
when paired --- $T$ exhibits the pair $\{1,3\}$. One $\size$ pass + \code{maxpart} scan
$=$ $\Theta(n)$ discovery (\coderef{c09-cent}, \thmref{c09-centfind}); the $\le n/2$
cap is exactly what makes recursing on pieces $\bigO{\log n}$-deep --- the decomposition
embryo Chapter~24 (\chref{c18}) finishes. Centroid $\ne$ centre (broom).
\selftest{
  \item List the pieces created by deleting a non-root node, with their sizes in
        $\size$-notation.
  \item Prove the ``at most one heavy piece'' fact and name the two theorems it powers.
  \item Why does equality ($\mathit{maxpart} = n/2$) exist, and what does it sit across?
  \item Give the broom's centroid and centre, and say which measurement finds each.
  \item What single condition gates CF~342E until Chapter~24?
}
\end{recallbox}
